\subsection{What transfers from the Diophantine and concise-iPRA projects}

The current \texttt{Computability/DiophantineConsistency/README.md}
was read at the pinned snapshot \cite{diophantine}.  Its constructions keep an exact witness
domain, exact polynomial or arithmetic circuit, and an exact native
refutation sentence tied together by model-uniform arguments and iPRA
endpoints.  That discipline transfers directly to the present use of exact
machine tables, exact source clocks, and the repository's native PA
consistency sentence.  Its small-variable and low-operation frontiers do
not by themselves provide a smaller binary Turing machine: an evaluator,
integer-witness enumeration, fixed-constant generation, and a blank startup
must all be realized in control states.  For example, the README's
compact44 family uses the fixed parameter $\beta=2^{590069}-1$.
A compact mathematical expression for that constant is not its cost in a
binary blank-start controller.  No claimed state saving is inferred from
the equation's variable count or operation count.

The current \texttt{Logic/PrimitiveRecursiveArithmetic/ConciseIPRA/README.md}
\cite{concise} measures complete standalone proof checkers by normalized Python tokens,
AST nodes, and bytes, and requires whole-domain bi-interpretability with
iPRA.  Its native scalar calculus uses iterator equations and a small
set of derived-logic rules.  These results guide the discipline of making
finite certificate checking and quantifier-free induction explicit, as in
the bounded-input compiler invariant above.  Substituting its proof checker
for the PA refutation search would, however, require a new calibration:
consistency of a theory bi-interpretable with iPRA is not automatically
consistency of PA.  Likewise, reducing Python syntax size is not a proven
reduction of literal binary Turing-machine controls.  The new constructions
therefore retain the PA-calibrated numerical calculus and use these
projects' proof-accounting methods, without claiming that either project
already supplies a smaller physical machine.
